\documentclass[11pt]{article}
\usepackage[T1]{fontenc}
\usepackage{lmodern}
\usepackage{microtype}
\usepackage{amsmath,amssymb,amsthm,mathtools}
\usepackage{booktabs}
\usepackage{geometry}
\usepackage{enumitem}
\usepackage{xcolor}
\usepackage{hyperref}
\geometry{margin=1in}
\hypersetup{
  colorlinks=true,
  linkcolor=blue!45!black,
  urlcolor=blue!55!black,
  pdftitle={Radical telescopes for Chebyshev transfers of every degree},
  pdfauthor={Magnus Gille}
}
\setlist{nosep,leftmargin=1.5em}
\setcounter{secnumdepth}{0}
\setlength{\parindent}{0pt}
\setlength{\parskip}{0.55em}
\emergencystretch=2em

\newtheorem{theorem}{Theorem}
\newtheorem{proposition}[theorem]{Proposition}
\newtheorem{corollary}[theorem]{Corollary}
\newtheorem{lemma}[theorem]{Lemma}
\theoremstyle{definition}
\newtheorem{definition}[theorem]{Definition}
\theoremstyle{remark}
\newtheorem{remark}[theorem]{Remark}

\newcommand{\rad}{\operatorname{rad}}
\newcommand{\Sq}{\operatorname{Sq}}
\newcommand{\T}{\mathcal T}
\newcommand{\Sset}{\mathcal S}

\title{Radical telescopes for Chebyshev transfers of every degree:\\
normalization, programmed squares, and the fixed-orbit boundary}
\author{Magnus Gille\\Independent researcher}
\date{Research draft, July 2026}

\begin{document}
\maketitle

\begin{quote}\small
This is a non-peer-reviewed companion research note.  It consolidates
results developed after the quadratic and prime-degree Chebyshev-orbit
manuscripts.  The transfer is classical and compositional.  The proposed
contribution is the degree-uniform arithmetic package described below;
specialist review of its proofs, sources, and priority remains necessary.
\end{quote}

\begin{abstract}
For every integer $d\ge2$, homogenizing $T_d$ gives an integral map
from a primitive positive triple $a+b=c$ of opposite parity to another
such triple.  We identify one fixed Lucas sequence behind the whole orbit.
After a necessary normalization at the primes dividing $d$, the quotient
between successive indices $d^j$ is a positive integer $E_{d,j}$.  Its
cyclotomic indices are precisely
\[
 \{m:m\mid d^{j+1},\ m\nmid d^j\}.
\]
For a natural class of $d$-admissible seeds, the $E_{d,j}$ are
pairwise coprime and disjoint from the seed.  This yields the exact
all-degree radical telescope
\[
 \frac{\rad(a_nb_nc_n)}{c_n}
 =\frac{\rad(a_0b_0c_0)}{c_0}
   \frac{|\sin(d^n\theta)|}
        {d^n\sin\theta\,W_{d,n}},
 \qquad
 W_{d,n}=\prod_{j<n}\frac{E_{d,j}}{\rad(E_{d,j})}.
\]
Composite degree bundles several Lucas atoms in one layer; prime degree is
the single-atom case.  We also complete the simultaneous finite-genealogy
realization for the quadratic transfer and prove a bounded iterated local
mean for the logarithmic defect in every prime degree, including $2$.
For one fixed orbit we isolate an unconditional growing congruence window
and show that every fixed valuation depth in a polynomial window is
harmless; the remaining obstruction is an exceptionally deep
Lucas--Wieferich tail.  A dual implementation checks that the canonical
quadratic, cubic, and quintic pairs have no new Lucas--Wieferich primes
between $10^6$ and $10^7$, and a 122-digit quadratic layer is completely
factored.  These results do not prove the fixed-orbit conjecture or any new
case of the $abc$ conjecture.

\medskip
\noindent\textit{2020 Mathematics Subject Classification.}
Primary 11B39; Secondary 11A05, 11D75, 11J86, 37P05.

\noindent\textit{Keywords.}
Chebyshev polynomial, Lucas sequence, Lucas atom, radical,
Wieferich prime, $abc$ conjecture, arithmetic dynamics.
\end{abstract}

\section{Introduction and novelty boundary}

Write
\[
 \rad(N)=\prod_{p\mid N}p.
\]
A primitive positive triple $a+b=c$ is an $abc$-hit when
\(\rad(abc)<c\).  Polynomial identities that transfer one additive triple
to another are classical; Martin and Miao survey many examples, and van der
Horst studies the transfer method in detail.  The Chebyshev identity used
here is a homogenized de Moivre identity, not a new polynomial identity.

Lucas atoms and their valuations have also been developed independently of
this project.  Sagan and Tirrell introduce the atom formalism, while Alecci,
Miska, Murru, and Romeo give a general definition and detailed $p$-adic
valuation results.  The finite-field factorization behind the local root
counts is equivalent, after a Cayley change of variable, to the fibotomic
factorization of Byer, Dvorachek, Eckard, Harrington, Wise, and Wong.
Ratliff, Rush, and Shah give related radical-preserving cyclotomic
constructions.

The proposed contribution is therefore narrower than ``a new Chebyshev
map.''  It is the following synthesis:

\begin{quote}\small
one normalization valid for every degree; prime-support separation of its
successive multi-atom layers; an exact radical telescope along the
resulting additive orbit; constructive control of arbitrary finite square
patterns in every prime degree; and a bounded local first moment whose
remaining integer-box obstruction is stated explicitly.
\end{quote}

Prime degree remains special: a new layer is a single cyclotomic atom (or
two coordinate branches of the same order), whereas composite degree
groups several atom indices.  Since $T_m\circ T_n=T_{mn}$, the
composite-degree transfer is itself only a block of lower-degree transfers.
The grouping and normalization, rather than the map, are the point.

The last sections distinguish finite-seed averages from a single fixed
orbit.  They sharpen the fixed-orbit obstruction but do not resolve it.
In particular, nothing here should be read as a proof of $abc$.

\section{The universal integral transfer}

Let $a+b=c$ be primitive and positive, with $a,b$ of opposite parity,
and put
\[
 D=b-a.
\]
Both $D$ and $c$ are odd.  Define the homogeneous Chebyshev values
\[
 H_r(D,c)=c^rT_r(D/c).
\]
They are integral because
\[
 H_0=1,\qquad H_1=D,\qquad
 H_{r+1}=2D H_r-c^2H_{r-1}.
\]

\begin{theorem}[Universal transfer]\label{thm:transfer}
For every integer $d\ge2$, set
\[
 \T_d(a,b,c)=
 \left(
 \frac{c^d-H_d(D,c)}2,
 \frac{c^d+H_d(D,c)}2,
 c^d
 \right).
 \tag{2.1}
\]
Then \(\T_d\) maps primitive positive opposite-parity triples to triples
of the same kind.  If
\[
 x=\frac{b-a}{c},
\]
then the new normalized difference is $T_d(x)$.  Moreover
\[
 \T_m\circ\T_n=\T_{mn}=\T_n\circ\T_m.
 \tag{2.2}
\]
After $n$ iterations,
\[
 c_n=c_0^{d^n},\qquad
 \frac{b_n-a_n}{c_n}
 =T_{d^n}\left(\frac{b_0-a_0}{c_0}\right).
 \tag{2.3}
\]
\end{theorem}

\begin{proof}
Since $H_d(D,c)$ is odd, the two coordinates in (2.1) are integers of
opposite parity and sum to $c^d$.  The inequality
$|T_d(x)|\le1$ gives nonnegativity.  Equality would make
$e^{i\arccos x}$ a root of unity, so $2D/c$ would be a rational
algebraic integer.  This is impossible because $c>1$, $c$ is odd, and
$\gcd(D,c)=1$.  Thus both coordinates are positive.

If an odd prime divided both new coordinates, it would divide $c$ and
$H_d(D,c)$.  But
\[
 H_d(D,c)\equiv 2^{d-1}D^d\pmod c,
\]
which is a unit modulo every prime divisor of $c$.  The prime $2$
cannot divide both coordinates because their sum is odd.  This proves
primitivity.  The difference of the two new coordinates is $H_d(D,c)$,
which proves the semiconjugacy.  Equations (2.2) and (2.3) follow from
$T_m\circ T_n=T_{mn}$.
\end{proof}

For $d=2$, (2.1) is exactly
\[
 (a,b,c)\longmapsto(4ab,(b-a)^2,c^2).
 \tag{2.4}
\]
For odd $d$, expanding
\((\sqrt b+i\sqrt a)^d\) splits the transfer into the familiar two
coordinate factors.  For even $d$, write instead
\[
 (\sqrt b+i\sqrt a)^d=C_d(a,b)+i\sqrt{ab}\,S_d(a,b).
\]
Then
\[
 \T_d(a,b)=
 \begin{cases}
 (aS_d(a,b)^2,\ bC_d(a,b)^2),&d\text{ odd},\\
 (abS_d(a,b)^2,\ C_d(a,b)^2),&d\text{ even}.
 \end{cases}
 \tag{2.5}
\]
Thus every degree has a two-coordinate-factor form, although the seed
prefactor depends on parity.

\section{One Lucas sequence and its normalized layers}

Fix the seed and put
\[
 \alpha=D+2\sqrt{-a_0b_0},\qquad
 \beta =D-2\sqrt{-a_0b_0}.
\]
Then
\[
 \alpha+\beta=2D,\qquad \alpha\beta=c_0^2,
 \qquad(\alpha-\beta)^2=-16a_0b_0.
\]
The associated Lucas sequence is
\[
 U_0=0,\quad U_1=1,\quad
 U_{r+1}=2D U_r-c_0^2U_{r-1},
 \qquad
 U_r=\frac{\alpha^r-\beta^r}{\alpha-\beta}.
 \tag{3.1}
\]

\begin{lemma}[Fixed-sequence identity]\label{lem:fixed-sequence}
For the degree-$d$ orbit,
\[
 a_nb_n=a_0b_0\,U_{d^n}^2.
 \tag{3.2}
\]
\end{lemma}

\begin{proof}
The norm identity
\(c_0^{2d^n}-(\alpha^{d^n}+\beta^{d^n})^2/4
=- (\alpha^{d^n}-\beta^{d^n})^2/4\), together with
\((\alpha-\beta)^2=-16a_0b_0\), gives (3.2).
\end{proof}

Define
\[
 Q_{d,j}=\frac{U_{d^{j+1}}}{U_{d^j}}.
 \tag{3.3}
\]
Homogeneous cyclotomic factorization gives
\[
 Q_{d,j}
 =\prod_{m\in\mathcal L_{d,j}}\Phi_m(\alpha,\beta),
 \qquad
 \mathcal L_{d,j}
 =\{m:m\mid d^{j+1},\ m\nmid d^j\}.
 \tag{3.4}
\]
The sets \(\mathcal L_{d,j}\) are pairwise disjoint.

The quotient (3.3) contains a systematic contribution from every prime
dividing $d$.  The following seed condition removes it exactly.

\begin{definition}
A seed is \emph{$d$-admissible} when it is primitive, positive, and of
opposite parity, and for every odd prime $q\mid d$,
\[
 q\mid a_0,\qquad v_q(U_q)=1.
 \tag{3.5}
\]
\end{definition}

For $q\ge5$, the second condition follows automatically from the first.
At $q=3$ it is necessary: the end terms in the binomial expression for
\(U_q/q\) can cancel modulo $q$.

\begin{theorem}[Normalized layer genealogy]\label{thm:genealogy}
Let $d\ge2$ and let the seed be $d$-admissible.  Then
\[
 E_{d,j}=\frac{|Q_{d,j}|}{d}
 \tag{3.6}
\]
is a positive integer coprime to $d a_0b_0c_0$.  The integers
\[
 E_{d,0},E_{d,1},E_{d,2},\ldots
\]
are pairwise coprime.  Consequently every prime outside the seed support
has at most one birth layer.
\end{theorem}

\begin{proof}
For an odd $q\mid d$, the Lucas law of repetition at a discriminant
prime, under (3.5), gives
\[
 v_q(U_r)=v_q(r)\qquad(q\mid r).
 \tag{3.7}
\]
Opposite parity gives the corresponding identity
\[
 v_2(U_r)=v_2(r)\qquad(2\mid r).
 \tag{3.8}
\]
These are standard specializations of Lucas-sequence valuation formulas;
the hypotheses include $q\nmid\alpha\beta$ and nondegeneracy.  Subtracting
the valuations at $d^{j+1}$ and $d^j$ gives
\[
 v_q(Q_{d,j})=v_q(d)\qquad(q\mid d),
\]
so (3.6) is integral and coprime to $d$.

If an odd $p\mid a_0b_0$, then
\(U_r\equiv r\alpha^{r-1}\pmod p\).  This is nonzero when $p\nmid d$,
and (3.7) removes the exact systematic part when $p\mid d$.  If
\(p\mid c_0\), one of \(\alpha,\beta\) is zero and the other is a unit
modulo $p$.  The parity assertion follows from (3.8) when $d$ is even
and directly from the recurrence at odd indices when $d$ is odd.  Hence
the normalized layers avoid the seed support.

Now suppose $p\nmid d a_0b_0c_0$ divided two layers.  At a prime above
$p$, a factor \(\Phi_m(\alpha,\beta)\) with $p\nmid m$ says that
\(\alpha/\beta\) has exact order $m$; inversion under conjugation does
not change that order.  The same $p$ cannot therefore occur at indices
from two disjoint sets \(\mathcal L_{d,j}\).  This proves pairwise
coprimality.
\end{proof}

\begin{remark}
For an odd prime $d=\ell$, a layer is the single prime-power
cyclotomic index \(\ell^{j+1}\), expressed through two coordinate
branches.  For $d=2$, (3.6) is \(|b_j-a_j|\).  A composite $d$ bundles
several atoms in (3.4).  Formula (3.6), not a termwise composite-binomial
argument, is the stable normalization.
\end{remark}

\section{The all-degree radical telescope}

Let
\[
 R_n=\rad(a_nb_nc_n),\qquad
 W_{d,n}=\prod_{j<n}\frac{E_{d,j}}{\rad(E_{d,j})}.
 \tag{4.1}
\]
Every prime divisor of $d$ already lies in the seed radical: the odd
ones divide $a_0$, and opposite parity handles $2$.  Telescoping (3.3)
and using Theorem~\ref{thm:genealogy} therefore gives
\[
 |U_{d^n}|=d^n\prod_{j<n}E_{d,j},\qquad
 R_n=R_0\prod_{j<n}\rad(E_{d,j}).
 \tag{4.2}
\]

Choose \(\theta\in(0,\pi)\) by
\[
 \cos\theta=\frac{b_0-a_0}{c_0}.
\]
Since \(\alpha=c_0e^{i\theta}\),
\[
 U_{d^n}=c_0^{d^n-1}
 \frac{\sin(d^n\theta)}{\sin\theta}.
 \tag{4.3}
\]

\begin{theorem}[All-degree radical telescope]\label{thm:telescope}
For every $d$-admissible seed and every $n\ge0$,
\[
 \boxed{
 \frac{R_n}{c_n}
 =\frac{R_0}{c_0}
  \frac{|\sin(d^n\theta)|}
       {d^n\sin\theta\,W_{d,n}}.}
 \tag{4.4}
\]
\end{theorem}

\begin{proof}
Divide the two identities in (4.2), substitute (4.3), and use
\(c_n=c_0^{d^n}\).
\end{proof}

For a fixed algebraic seed, a standard linear-form estimate supplies a
constant \(C=C(a_0,b_0,d)\) such that
\(\log|\sin(d^n\theta)|\ge-C(n+1)\);
the elementary upper bound is immediate.  Thus the exponential-scale
obstruction in (4.4) is precisely the powerful part $W_{d,n}$.  In the
quadratic and prime-degree cases this yields the familiar equivalence
\[
 \log W_{d,n}=o(d^n)
 \quad\Longleftrightarrow\quad
 \text{the orbit quality tends to }1.
 \tag{4.5}
\]
A persistent quality gap would require \(\log W_{d,n}\) to have positive
scale $d^n$.  Equation (4.4) is an exact reduction, not a bound on that
powerful part.

\section{Programmed squares in prime degree}

The prime-degree manuscript proves that any finite compatible collection
of birth levels, branches, splitting signs, and exact multiplicities can
be realized simultaneously for every odd prime degree.  Here we record
the missing quadratic counterpart.  It is also useful as a source of exact
test vectors for the fixed-orbit discussion.

For the quadratic transfer (2.4), let
\[
 G_n(X,Y)=b_n-a_n
\]
be the homogeneous seed-coordinate form and put \(m=2^{n+2}\).

\begin{proposition}[Quadratic local roots]\label{prop:quadratic-roots}
Let $p$ be odd.  If $m\mid p-\chi$ for
\(\chi\in\{1,-1\}\), then \(G_n(X,1)\) has exactly $2^n$ distinct
simple roots in \(\mathbb F_p\).  They are
\[
 \rho_\zeta=-\left(\frac{\zeta-1}{\zeta+1}\right)^2,
 \tag{5.1}
\]
where \(\zeta\) has exact order $m$ in \(\mathbb F_p^\times\) when
\(\chi=1\), or in the norm-one subgroup of \(\mathbb F_{p^2}^\times\)
when \(\chi=-1\), modulo \(\zeta\sim\zeta^{-1}\).  There are no roots
when neither divisibility holds.  Every root satisfies
\[
 \rho\ne0,-1,\qquad
 \left(\frac{-\rho}{p}\right)=\chi.
 \tag{5.2}
\]
Each root has a unique Hensel lift, and if \(\widehat\rho\) is its lift
modulo $p^h$, then for every unit \(\lambda\),
\[
 v_p\bigl(G_n(\widehat\rho+\lambda p^h,1)\bigr)=h.
 \tag{5.3}
\]
\end{proposition}

\begin{proof}
In the quadratic extension generated by square roots of $X$ and $Y$,
the standard Cayley parameter identifies the roots of $G_n$ with exact
order-$m$ roots of unity, and inversion has the same seed-coordinate
image.  The split and norm-one groups have orders $p-1$ and $p+1$,
respectively.  Degree exhaustion gives \(\varphi(m)/2=2^n\) roots.
The exceptional values $0,-1$ are not roots, and the derivative does not
vanish because $p\nmid m$.  Formula (5.2) follows directly from (5.1).
Hensel's lemma and the first nonzero Taylor term give (5.3).
\end{proof}

\begin{theorem}[Simultaneous quadratic realization]\label{thm:quadratic-realization}
Let
\[
 \mathcal D=\{(p_i,n_i,h_i,\chi_i):1\le i\le r\}
\]
have distinct odd primes, $n_i\ge0$, $h_i\ge1$, and
\[
 2^{n_i+2}\mid p_i-\chi_i.
\]
There are infinitely many primitive positive seeds with $a_0$ even and
$b_0$ odd such that, simultaneously for every $i$,
\[
 v_{p_i}(E_{2,n_i})=h_i,
 \qquad
 \left(\frac{D_K}{p_i}\right)=\chi_i,
 \tag{5.4}
\]
where \(K=\mathbb Q(\sqrt{-a_0b_0})\), while $p_i$ divides neither the
seed nor any layer $E_{2,j}$ with $j\ne n_i$.
\end{theorem}

\begin{proof}
For each datum choose a root from Proposition~\ref{prop:quadratic-roots}
and prescribe its Hensel lift with the perturbation in (5.3) modulo
\(p_i^{h_i+1}\).  Fix $b_0=1$.  The Chinese remainder theorem, together
with \(a_0\equiv0\pmod2\), supplies infinitely many positive $a_0$ in
the resulting progression.  These seeds are primitive and have the
required parity.  Equations (5.2) and (5.3) give the splitting sign and
exact valuation.  The exclusions \(\rho\ne0,-1\) keep $p_i$ away from
\(a_0b_0c_0\).  Finally, the Cayley parameter has exact order
\(2^{n_i+2}\), so it cannot belong to the distinct order class attached to
another layer.  This proves the last assertion.
\end{proof}

Theorem~\ref{thm:quadratic-realization}, together with the two-branch
odd-prime result, shows that arbitrarily deep but finitely prescribed
square behavior is constructible.  It does not produce a fixed seed with
large defect at infinitely many layers.

\section{A bounded local mean for every prime degree}

This section concerns seed averages, not one orbit.  Let $d$ be prime.
For $d=2$, let \(\Sset_2(H)\) be the primitive opposite-parity pairs in
\([1,H]^2\).  Its density is
\[
 \kappa_2=\frac{2}{3\zeta(2)}.
 \tag{6.1}
\]
For an odd prime $d=\ell$, take the admissible seed set of the
prime-degree manuscript.  Its density is
\[
 \kappa_\ell=
 \begin{cases}
 1/(9\zeta(2)),&\ell=3,\\
 2/(3(\ell+1)\zeta(2)),&\ell\ge5.
 \end{cases}
 \tag{6.2}
\]

Put
\[
 q_j=
 \begin{cases}
 2^{j+2},&d=2,\\
 d^{j+1},&d\text{ odd},
 \end{cases}
 \qquad
 \nu_d(j)=
 \begin{cases}
 \varphi(q_j)/2=2^j,&d=2,\\
 \varphi(q_j)=d^j(d-1),&d\text{ odd}.
 \end{cases}
 \tag{6.3}
\]
The factor $1/2$ in the quadratic case comes from inversion in its
single Cayley tower; the odd-prime count combines two branches.

For fixed $n,P,K$, define the bounded defect
\[
 D_{n,P,K}(a,b)=
 \sum_{j<n}\ \sum_{\substack{p\le P\\p\nmid2d}}
 \log p\sum_{h=2}^{K+1}\mathbf 1_{p^h\mid E_{d,j}(a,b)}.
 \tag{6.4}
\]

\begin{theorem}[Bounded iterated local mean]\label{thm:local-mean}
For every prime degree $d\ge2$,
\[
 \lim_{P\to\infty}\lim_{K\to\infty}\lim_{H\to\infty}
 \frac{1}{|\Sset_d(H)|}
 \sum_{(a,b)\in\Sset_d(H)}D_{n,P,K}(a,b)
 =L_d(n),
 \tag{6.5}
\]
where
\[
 L_d(n)=
 \sum_{j<n}\nu_d(j)
 \sum_{\substack{p\nmid2d\\q_j\mid p-1\ \mathrm{or}\ q_j\mid p+1}}
 \frac{\log p}{p^2-1}.
 \tag{6.6}
\]
There is an effective constant $C_d$, independent of $n$, such that
\[
 L_d(n)\le C_d<\infty.
 \tag{6.7}
\]
For odd prime $d$, one may take
\[
 C_d=\frac58\left[
 \frac{\zeta(2)\log3-\zeta'(2)}{d-1}
 +\frac{\zeta(2)d\log d}{(d-1)^2}
 \right].
 \tag{6.8}
\]
For $d=2$, one may take
\[
 C_2=\frac32\zeta(2)\log2-\frac38\zeta'(2).
 \tag{6.9}
\]
In particular, \(C_d\ll\log d/(d-1)\) along odd prime degrees.
\end{theorem}

\begin{proof}
At a compatible prime, the relevant level form has exactly
\(\nu_d(j)\) simple projective roots.  Each has one lift modulo $p^h$
and \(\varphi(p^h)\) unit representatives.  Among the
\(p^{2h}-p^{2h-2}\) primitive pairs modulo $p^h$, this gives
\[
 \Pr(p^h\mid E_{d,j})
 =\frac{\nu_d(j)}{p^{h-1}(p+1)}.
 \tag{6.10}
\]
The conditions at $2$ and $d$ separate by the Chinese remainder
theorem.  Summing (6.10) for \(2\le h\le K+1\) gives
\[
 \nu_d(j)\frac{1-p^{-K}}{p^2-1}.
\]
For fixed cutoffs the box limit is ordinary counting in finitely many
residue classes, which proves (6.5) before the last two limits.

For odd $q\ge3$, the eligible integers are \(2rq\pm1\), and
\[
 \sum_{q\mid p^2-1}\frac{\log p}{p^2-1}
 \le\frac5{8q^2}
 \left(\zeta(2)\log(3q)-\zeta'(2)\right).
 \tag{6.11}
\]
Multiplying by \(\varphi(d^{j+1})\) and summing the two geometric series
gives (6.8).  When $q=2^{j+2}$, writing eligible integers as
\(rq\pm1\) gives the elementary majorant
\[
 \sum_{q\mid p^2-1}\frac{\log p}{p^2-1}
 \le\frac3{q^2}
 \left(\zeta(2)\log(2q)-\zeta'(2)\right).
 \tag{6.12}
\]
Since \(\nu_2(j)=q/4\), summation yields (6.9).
\end{proof}

The theorem deliberately uses the order of limits displayed in (6.5).
Passing from this local model to the untruncated integer-box mean would
require uniform control of primes growing with the box:
\[
 \lim_{P\to\infty}\limsup_{H\to\infty}
 \mathbb E_{\Sset_d(H)}
 \left[\log W_{d,n}-D_{n,P,\infty}\right]=0.
 \tag{6.13}
\]
This is a large-square-divisor problem for binary forms.  Fixed-modulus
Hensel densities do not prove (6.13), and general multivariable results
that assume $abc$ cannot be used here without circularity.

For composite degrees, (3.4) suggests the convergent majorant
\[
 \sum_{\substack{m\mid d^\infty\\m>1}}
 \frac{\varphi(m)\log(3m)}{m^2}.
\]
A theorem would require a fully checked simultaneous root classification
for every composite-layer atom.  We leave that extension open rather than
promote the majorant to a result.

\section{What remains for one fixed orbit}

Restrict now to the quadratic orbit or one admissible odd-prime-degree
orbit.  Write $E_j=E_{d,j}$,
\[
 \delta_j=\log\frac{E_j}{\rad(E_j)},\qquad
 q_j=\begin{cases}2^{j+2},&d=2,\\d^{j+1},&d\text{ odd}.
 \end{cases}
 \tag{7.1}
\]
Every sufficiently large prime $p\mid E_j$ satisfies
\[
 p\equiv\pm1\pmod{q_j}.
 \tag{7.2}
\]

Let
\[
 \Sq(N)=\prod_p p^{\lfloor v_p(N)/2\rfloor}.
\]
Nonnegativity and
\[
 \log\Sq(E_j)\le\delta_j\le2\log\Sq(E_j)
\]
give the exact equivalence
\[
 \log W_{d,n}=o(d^n)
 \Longleftrightarrow
 \delta_j=o(d^j)
 \Longleftrightarrow
 \log\Sq(E_j)=o(d^j).
 \tag{7.3}
\]
Thus the fixed-orbit problem is a subpower bound for the largest square
dividing one sparse sequence of cyclotomic values.

There is one unconditional growing window.  Stewart's valuation lemma for
a fixed Lucas pair implies, outside an effective finite set,
\[
 \operatorname{ord}_{\mathfrak p}(u^m-1)
 <p\exp\left(-\frac{\log p}{51.9\log\log p}\right)
 \log|\alpha|\log m,
 \tag{7.4}
\]
where $u=\alpha/\beta$.

\begin{theorem}[A growing harmless window]\label{thm:stewart-window}
Put \(L_j=\log q_j/\log\log q_j\).  For every fixed
\(0<\gamma<1/103.8\),
\[
 \sum_{\substack{p\le q_j\exp(\gamma L_j)\\p^2\mid E_j}}
 (v_p(E_j)-1)\log p=o(q_j).
 \tag{7.5}
\]
\end{theorem}

\begin{proof}
Choose $a'$ with \(2\gamma<a'<1/51.9\).  Since every candidate prime
is at least $q_j-1$, (7.4) contributes a factor at most
\(e^{-a'L_j}\).  There are at most
\(2(e^{\gamma L_j}+1)\) candidate integers in (7.2).  Multiplying the
number of candidates, the window endpoint, the valuation bound, and
\(\log p\) gives
\[
 q_j\exp((2\gamma-a'+o(1))L_j)=o(q_j).
\]
\end{proof}

This is larger than every fixed multiple of $q_j$, but smaller than
\(q_j^{1+\varepsilon}\) for each fixed \(\varepsilon>0\).  Inside such a
polynomial window, the exact remaining obstruction can be stated
elementarily.  For $Y\ge q_j$, put
\[
 D_j(Y)=\sum_{\substack{p\le Y\\p^2\mid E_j}}
 (v_p(E_j)-1)\log p.
\]
For an integer $K\ge1$, split it as
\[
 \begin{aligned}
 T_{j,K}(Y)&=\sum_{\substack{p\le Y\\p^2\mid E_j}}
 \min\{v_p(E_j)-1,K\}\log p,\\
 R_{j,K}(Y)&=\sum_{\substack{p\le Y\\p^2\mid E_j}}
 (v_p(E_j)-1-K)_+\log p.
 \end{aligned}
 \tag{7.6}
\]

\begin{theorem}[Polynomial window equals its deep tail]\label{thm:deep-tail}
Let \(1\le\beta<2\), \(Y=q_j^\beta\), and
\[
 K_j=o\left(\frac{q_j^{2-\beta}}{\log q_j}\right).
\]
Then
\[
 T_{j,K_j}(Y)=o(q_j),\qquad
 D_j(Y)=o(q_j)\Longleftrightarrow R_{j,K_j}(Y)=o(q_j).
 \tag{7.7}
\]
\end{theorem}

\begin{proof}
By (7.2), at most \(2(Y/q_j+1)\) candidate integers occur.  Therefore
\[
 T_{j,K_j}(Y)
 \le2K_j(Y/q_j+1)\log Y=o(q_j).
\]
The equivalence follows from the nonnegative exact split (7.6).
\end{proof}

In particular, squares, cubes, every fixed valuation depth, and depths
growing almost as quickly as \(q_j^{2-\beta}/\log q_j\) are harmless.
The problem is not merely to show that Lucas--Wieferich primes are rare;
it is to rule out an aggregate of extraordinarily deep lifts.

This phrasing is exact.  Away from the seed and ramified primes,
\(p^2\mid E_j\) forces
\[
 v_p\bigl(U_{p-(\Delta/p)}\bigr)=v_p(E_j)\ge2.
 \tag{7.8}
\]
Thus every squared layer prime is Lucas--Wieferich for the fixed pair.
If its valuation is at least $3$, the corresponding prime ideal is
super-Wieferich.  Finiteness of those super-Wieferich prime ideals would,
for example, imply \(D_j(q_j^\beta)=o(q_j)\) for every fixed
\(\beta<2\); that finiteness is an open hypothesis, not an available
theorem.

\section{Deeper finite computation}

Finite computation neither proves nor predicts the fixed-orbit
asymptotic.  It does give exact regression targets for the layer formulas
and the Lucas--Wieferich bridge.

\subsection{The 122-digit quadratic layer}

For the quadratic seed \((a_0,b_0,c_0)=(1,8,9)\), the level-seven
normalized layer is the 122-digit integer
\begin{align*}
 E_7={}&7899993675270888986790467637030360278248237512403075459722095\\
      &7868119418164413649973647970771585795292694673449586011206657.
\end{align*}
GMP-ECM 7.0.7 split the formerly unresolved 98-digit residual by one
Pollard $p-1$ step and two ECM discoveries.  Exact division gives
\begin{align*}
E_7={}&189439\cdot750692351\cdot9825841153\cdot298196593663\\
&\cdot991245449894911\cdot6726631000961507661177857\\
&\cdot28434404151626641091139435909034910237447173121.
\tag{8.1}
\end{align*}
All seven factors pass independent primality tests, their exact product is
recomputed from the recurrence, each occurs to exponent one, and their
residues modulo $512$ are
\[
 511,\ 511,\ 1,\ 511,\ 511,\ 1,\ 1.
\]
Thus the forced \(p\equiv\pm1\pmod{512}\) condition holds and this layer
has zero repeated-prime defect.

\subsection{A census beyond one million}

For a Lucas pair \((P,Q)\), the census tests every odd unramified prime
\(p\le10^7\) for
\[
 p^2\mid U_{p-(\Delta/p)}(P,Q)
\]
and then tests the hits modulo $p^3$.  Two implementations were used:
binary exponentiation in \(\mathbb Z[X]/(X^2-PX+Q)\), and an independently
written companion-matrix algorithm.  Each scanned all 664,578 odd primes
through \(10^7\), and the outputs agree exactly.

\begin{center}
\begin{tabular}{@{}lrrrr@{}}
\toprule
orbit pair \((P,Q)\) & Wieferich \(p\) & rank & tower rank? & depth \(\ge3\)?\\
\midrule
quadratic \((14,81)\) & 65519 & 455 & no & no\\
cubic \((-2,25)\) & 47 & 24 & no & no\\
quintic \((-6,49)\) & 53 & 26 & no & no\\
\bottomrule
\end{tabular}
\end{center}

There are no new hits between \(10^6\) and \(10^7\).  None of the three
known hits can enter its displayed orbit tower because its rank is not a
pure power of the orbit degree, and none is super-Wieferich.  This is a
bounded observation only; no independence or density assertion is being
made.

The exact inputs, intermediate cofactors, software versions, machine-readable
census output, and regression commands accompany the source repository.

\section{Conclusions}

The all-degree result is structural.  It shows that the quadratic orbit,
the odd-prime construction, and every composite block share one normalized
Lucas-layer mechanism and one radical identity.  The quadratic realization
and the prime-degree local mean complete two genuine gaps between the first
two manuscripts.

The fixed-orbit result is more modest.  Congruence and valuation theory
control a growing \(q_j^{1+o(1)}\) window, and pure counting disposes of
every bounded or moderately growing lift depth in any fixed polynomial
window below \(q_j^2\).  What survives is a deep diagonal
Lucas--Wieferich tail.  Neither the complete level-seven factorization nor
the census through \(10^7\) changes that logical boundary.

Accordingly, these papers have moved closer to $abc$ in the sense of
producing an exact reduction and a sharper partial bound, not in the sense
of proving a new case of the conjecture.  A next theorem would need genuine
uniform control of deep lifts for one fixed Lucas pair, or a different
global mechanism for the powerful part.

\section{AI-use statement}

OpenAI Codex and Anthropic Claude were used during Phases 5--8 for
mathematical exploration, proof and source audits, independent software
implementations, computation, drafting, and adversarial cross-review.  The
named author must independently verify every theorem, proof, reference,
novelty statement, and computation before circulation or submission.  The
AI systems are not authors.

\section{References}

G. Alecci, P. Miska, N. Murru, and G. Romeo, ``On alternative definition
of Lucas atoms and their \(p\)-adic valuations,'' \emph{Monatshefte f\"ur
Mathematik} \textbf{207} (2025), 175--196.
\url{https://doi.org/10.1007/s00605-025-02087-w}

C. Byer, T. Dvorachek, E. Eckard, J. Harrington, L. Wise, and
T. W. H. Wong, ``On the properties of fibotomic polynomials,''
\emph{Advances in Applied Mathematics} \textbf{138} (2022), 102344.
\url{https://doi.org/10.1016/j.aam.2022.102344}

G. Martin and W. Miao, ``\(abc\) triples,'' \emph{Functiones et
Approximatio Commentarii Mathematici} \textbf{55} (2016), 145--176.
\url{https://doi.org/10.7169/facm/2016.55.2.2}

B. Poonen, ``Squarefree values of multivariable polynomials,''
\emph{Duke Mathematical Journal} \textbf{118} (2003), 353--373.
\url{https://doi.org/10.1215/S0012-7094-03-11826-8}

L. J. Ratliff, Jr., D. E. Rush, and K. Shah, ``Note on cyclotomic
polynomials and prime ideals,'' \emph{Communications in Algebra}
\textbf{32} (2004), 333--343.
\url{https://doi.org/10.1081/AGB-120027870}

B. E. Sagan and J. P. Tirrell, ``Lucas atoms,'' \emph{Advances in
Mathematics} \textbf{374} (2020), 107387.
\url{https://doi.org/10.1016/j.aim.2020.107387}

C. Sanna, ``The \(p\)-adic valuation of Lucas sequences,''
\emph{The Fibonacci Quarterly} \textbf{54} (2016), 118--124.
\url{https://www.fq.math.ca/Papers1/54-2/Sanna02242016.pdf}

M. Gille, ``Radicals in iterated quadratic \(abc\)-transfers,'' research
draft, 2026.

M. Gille, ``Prime genealogy in Chebyshev \(abc\)-orbits: interleaved
Lucas atoms and exact radical telescopes,'' research draft, 2026.

C. L. Stewart, ``On divisors of Lucas and Lehmer numbers,''
\emph{Acta Mathematica} \textbf{211} (2013), 291--314.
\url{https://doi.org/10.1007/s11511-013-0103-2}

J. P. van der Horst, \emph{Finding \(ABC\)-triples using elliptic curves},
M.Sc. thesis, Universiteit Leiden, 2010.
\url{https://math.leidenuniv.nl/scripties/vanderHorstMaster.pdf}

\end{document}
